% Candidate strengthening of Notebook 12; requires independent audit.
% Prefix ss-. Depends on rd heat estimates and flat alias notation.
\section{Equal-exponent Schauder bounds and the scalar inverse}
\label{ss-section}

The positive H\"older loss in the preceding right-divergence
theorem is sufficient for the displayed source estimates. It is
not necessary for a fixed exponent strictly between zero and one.
The following refinement is useful for nonlinear elliptic equations,
where an exponent loss at every iteration would not close a fixed
point argument.

\begin{theorem}\label{ss-equal-exponent}
For every fixed $0<\alpha<1$, the homogeneous flat operators obey
\begin{equation}\label{ss-main}
 N_{\alpha,\varpi}(GD^*\lambda)
       \le C_\alpha\norm{\lambda}_{\alpha,\varpi},
 \qquad G\in\{G_0,G_\pi^J,G_1^J\}.
\end{equation}
The spatial exponent is sufficiently small, and the current is in
the common ball for the $N$-to-$N_\alpha$ estimates established
above. Constants are uniform in refinement and compatible volume.
There is no uniform assertion as $\alpha$ tends to zero or one.
\end{theorem}

\begin{proof}
Use $G_0=G_m+R_m$ from \eqref{rd-massive-split}.
Equation \eqref{rd-remainder-bound} already controls $R_mD^*$
from bounded input into $N_\alpha$ for every $\alpha<1$.
Only the gradient H\"older member of $G_mD^*$ needs a new
argument. The heat moment and exact row cancellation, with
$\gamma=\alpha$, give
\[
 \norm{D^2e^{-t\mathcal L}\lambda}_{\infty,\varpi}
       \le Ct^{-1+\alpha/2}\norm{\lambda}_{\alpha,\varpi},
 \qquad
 \norm{D^3e^{-t\mathcal L}\lambda}_{\infty,\varpi}
       \le Ct^{-3/2+\alpha/2}\norm{\lambda}_{\alpha,\varpi}
\]
for $0<t\le1$, including times below the mesh scale.
For two fine output sites in a fixed collar at distance
$0<r\le1$, estimate their second-derivative difference by the
first bound for $t<r^2$, and by $r$ times the second for $t\ge r^2$.
The finite path used for this difference stays in a fixed enlarged
collar, so its weight comparison costs only a fixed constant.
Consequently the difference after integrating the massive heat
kernel is bounded, divided by the input norm, by
\begin{align}
 C\int_0^{r^2}t^{-1+\alpha/2}\,dt
 +Cr\int_{r^2}^1t^{-3/2+\alpha/2}\,dt+Cr
 &\le C\left(\frac2\alpha+
                     \frac2{1-\alpha}+1\right)r^\alpha.
 \label{ss-split-integral}
\end{align}
The last term is the large-time massive tail, using its third
derivative bound. Distances above one are controlled by the
supremum bound. The latter uses
$\int_0^1t^{-1+\alpha/2}dt=2/\alpha$ for the gradient and the
integrable first heat derivative for the field. This proves the
baseline assertion. It improves the preceding argument by
splitting the time integral for each output separation before
taking the H\"older supremum, rather than integrating a fixed
H\"older operator norm down to zero time.

The current transfer is exactly \eqref{rd-current-transfer}:
the baseline datum is in $N_\alpha$, and the perturbation maps
$N$ to $N_\alpha$ with a strict contraction in $N$. Its radius
does not depend on the initial datum constant. Periodic extension,
complete weighted fields and smooth localization work as before.
\end{proof}

\subsection{A separately normalized scalar statement}

Let
\[
 G'=(\mathcal L+Q_0^*a'Q_0)^{-1},\qquad
              0<a'_-\le a'\le a'_+<\infty,
\]
where the distinct scalar penalty ranges in a fixed compact
positive interval. This hypothesis is explicit: no unprinted
scale-dependent scalar penalty is assigned a value here.
With the scalar aliases $a_r$, put
\begin{equation}\label{ss-scalar-denominator}
 A=a_0a_0^-,\quad
 R_{s,1}=\sum_{r\ne0}\frac{a_ra_r^-}{P_r},\quad
 \widehat R_{s,1}=a'^{-1}+R_{s,1},\quad
 D_{s,1}=A+P_0\widehat R_{s,1}.
\end{equation}
The canceled scalar Woodbury formula is
\begin{align}
 G'_{00}&=\widehat R_{s,1}/D_{s,1},\nonumber\\
 G'_{0s}&=-a_0^-a_s/(P_sD_{s,1}),
 &G'_{r0}&=-a_r^-a_0/(P_rD_{s,1}),\nonumber\\
 G'_{rs}&=\frac{\delta_{rs}}{P_r}
       -\frac{P_0a_r^-a_s}{P_rP_sD_{s,1}},
 &&r,s\ne0.
 \label{ss-scalar-formula}
\end{align}
These formulas extend across $P_0=0$; no singular factor $1/P_0$
remains. They follow by the rank-one inverse identity away from
$P_0=0$ and analytic continuation. On the real compact chart,
$D_{s,1}\ge A\ge c>0$. Alias Parseval bounds and the high-alias
gap give uniform derivative bounds, so $D_{s,1}$ remains bounded
away from zero on a common complex strip. The constants may
depend on $a'_-,a'_+$.

Subtract the scalar massive inverse $G_m=(\mathcal L+1)^{-1}$.
The high diagonal has exactly the improved heat representation
\eqref{rd-high-remainder}. All other terms in
\eqref{ss-scalar-formula} have bounded central coefficients and
high output vector $a_r^-/P_r$. This vector has uniform absolute
alias summability after $1+\alpha$ physical differences when
$\alpha<1$: the bound
$|a_r|\le C\prod_\nu(1+|r_\nu|)^{-1}$ gives a dyadic shell sum
bounded by $C(1+n)^{d-1}2^{-n(1-\alpha)}$. One right derivative
leaves the input row bounded in alias $\ell^2$, since
$|u_ra_r/P_r|\le C|a_r|/(1+|r|)$ for high aliases.
The undifferentiated input rows have the same property.
Uniform cell Fourier conversion, periodic gluing and contour
shifting therefore give
\begin{equation}\label{ss-scalar-remainder}
 N_{\alpha,\varpi}\bigl((G'-G_m)D^*f\bigr)
       \le C_\alpha\norm{f}_{\infty,\varpi}.
\end{equation}
The bounded central massive subtraction is included in this
estimate. Only the complete reconstructed operator is periodized
and shifted; the individual central pieces need not be periodic.

Applying the same massive cancellation proof establishes
\begin{equation}\label{ss-scalar-right-divergence}
 N_{\alpha,\varpi}(G'D^*\lambda)
       \le C_\alpha\norm{\lambda}_{\alpha,\varpi},
                       \qquad0<\alpha<1.
\end{equation}
Using instead the preceding positive-loss proof also gives
$N_{\beta,\varpi}(G'D^*\lambda)
\le C_{\beta,\gamma}\norm{\lambda}_{\gamma,\varpi}$ for
$0\le\beta<\gamma<1$. This proves the scalar analytic estimate
for the stated penalty normalization. Application to a separately
specified source penalty requires verifying its uniform positive
interval. It neither replaces $G'$ by the one-form $G_0$ nor
identifies their gauge projections.

Equal H\"older exponents strictly between zero and one in
\eqref{ss-main} do not assert a Lipschitz-gradient endpoint or a
second-difference bound on arbitrary bounded forcing. The earlier
endpoint counterexample remains applicable to those other claims.
